\clearpage
\fancyhead[L]{\small CERES.}
\begin{center}
CERES.

\rule{0.6\linewidth}{0.4pt}

Monatliche Correspondenz. Band XI. S. 283–292. März 1805.

\rule{0.6\linewidth}{0.4pt}
\end{center}

{\rightskip=0pt plus 1em
Auch die Zeit, in der Ceres sichtbar war, ist nun vorbei. Selbst mit größeren Äquatorialinstrumenten dürfte ihre Beobachtung schwierig sein, weil sie in beträchtlicher Entfernung vom Meridian stattfinden müsste. Wir müssen nun sechs bis sieben Monate darauf verzichten, diesen Wandelstern am Himmel zu verfolgen, und können nur darauf hinarbeiten, ihn bei seinem nächsten Erscheinen sofort wiederzufinden. Mit den von Dr. \textsc{Gauss} erneut verbesserten Elementen und der daraus berechneten Ephemeride wird das keine Schwierigkeiten bereiten, da man die Bahn der Ceres jetzt als sehr genau bestimmt ansehen kann. Wir geben hier noch alles an, was ihre künftige Beobachtung erleichtern kann.\par}

Die nach den IX. Elementen berechnete Ephemeride des geozentrischen Laufs der Ceres gab ihre Rektaszension im September und Oktober um neun Minuten zu klein und ihre Deklination um vier Minuten zu groß an. Deshalb zog Dr. \textsc{Gauss} die drei Oppositionen der Ceres in den Jahren 1802, 1803 und 1804 heran und stützte darauf die folgenden X. Elemente.
\SourceLeaf{267}
\begin{center}
\begingroup\fontsize{9}{11}\selectfont\setlength{\tabcolsep}{3pt}\renewcommand{\arraystretch}{1.1}
\begin{tabular}{lr}
Epoche 1804, Seeberger Meridian . . . . . . . . & $312^\circ\ 1'\ 33''5$ \\
Tägliche Bewegung . . . . . . . . . . . . . & $771''0524$ \\
Aphel . . . . . . . . . . . . . . . & $326^\circ\ 26'\ 3''1$ \\
Exzentrizität . . . . . . . . . . . . . . . & $0.0784757$ \\
Logarithmus der Halbachse . . . . . . . . . & $0.04420004$ \\
$\AscendingNode$ 1804 . . . . . . . . . . . . . . . . . & $80^\circ\ 59'\ 12''$ \\
Neigung . . . . . . . . . . . . . . . . . & $10\ 37\ 45$ \\
\end{tabular}\endgroup
\end{center}

\indent Aus diesen Elementen berechnete Dr. \textsc{Gauss} den geozentrischen Lauf dieses Planeten vom 28. Juli 1805 bis zum 24. Mai 1806 [und gibt die Rektaszension und die Deklination minutengenau an, den Abstand von der Erde auf drei Dezimalstellen und die Helligkeit auf fünf Dezimalstellen, jeweils um 12 Uhr Seeberger Zeit in Abständen von drei Tagen.]

—\quad—\quad—

Er [\textsc{Gauss}] schrieb uns dazu, dass \textsc{Bessel}, der durch seine Abhandlung über die Bahn des Kometen von 1607 (M. C. 1804, Bd. X, S. 425 f.) zu großem Ansehen gelangt war, ihm bei der Berechnung dieser Ephemeride half, indem er alle dafür nötigen Sonnenpositionen lieferte. Wer die Rechnungen kennt, die zur Bestimmung der Elemente eines Planeten und dann zur Ableitung jeder einzelnen Position daraus nötig sind, muss gewiss bewundern, wie ein einzelner Mann in so kurzer Zeit so viele mühsame Rechnungen vollenden konnte. —\quad—\quad—

\hfill[\textsc{von Zach.}]
